\chapter{Geração de Código WebAssembly}\label{chap:wat}

\section{Introdução}

Nos capítulos anteriores construímos a \emph{representação intermediária em
árvore} (IRT) e aplicamos otimizações sobre ela.  O passo final do compilador é
a \textbf{seleção de instruções}: traduzir a IRT para uma linguagem executável
pela máquina-alvo.

Neste capítulo a máquina-alvo é o \textbf{WebAssembly} (Wasm), uma
especificação de máquina virtual com semântica formal publicada pelo W3C.
WebAssembly é uma escolha excelente para fins pedagógicos: sua semântica é
precisa, sua representação textual (WAT) é legível, e ferramentas como
\texttt{wabt} e \texttt{wasmtime} permitem validar e executar o código gerado
diretamente a partir da linha de comando.

\begin{center}
  \begin{tikzpicture}[
      node distance=1.2cm and 1.5cm,
      box/.style={draw, rounded corners, minimum width=2.8cm, minimum
        height=0.9cm, font=\small\sffamily, align=center},
      arr/.style={-Stealth, thick}
    ]
    \node[box] (irt)  {IRT};
    \node[box, right=of irt]  (opt)   {IRT otimizada};
    \node[box, right=of opt]  (wat)   {Módulo WAT};
    \node[box, right=of wat]  (wasm)  {Binário Wasm};
    \draw[arr] (irt) -- (opt) node[midway, above, font=\tiny\sffamily]{opt.};
    \draw[arr] (opt) -- (wat) node[midway, above, font=\tiny\sffamily]{codegen};
    \draw[arr] (wat) -- (wasm) node[midway, above, font=\tiny\sffamily]{\texttt{wat2wasm}};
  \end{tikzpicture}
\end{center}

O gerador descrito neste capítulo recebe um programa IRT (possivelmente
otimizado) e produz um módulo WAT que pode ser montado com \texttt{wat2wasm} e
executado com \texttt{wasmtime}.

% ============================================================================
\section{O Formato Textual WebAssembly (WAT)}\label{sec:wat-syntax}
% ============================================================================

\subsection{Visão Geral}

Um \textbf{módulo} WAT é um arquivo de texto estruturado em S-expressões.  Cada
módulo agrupa declarações de importações, memória, variáveis globais, funções e
exportações.  A figura a seguir mostra a estrutura geral:

\begin{minted}{text}
(module
  (import "wasi_snapshot_preview1" "fd_write"
          (func $fd_write (param i32 i32 i32 i32) (result i32)))
  (import "wasi_snapshot_preview1" "fd_read"
          (func $fd_read  (param i32 i32 i32 i32) (result i32)))
  (memory (export "memory") 4)
  (global $heap_ptr (mut i32) (i32.const 128))
  (func $alloc  ...)
  (func $main   ...)   ;; funções do usuário
  (func $print  ...)   ;; runtime I/O embutido
  (func $read_int ...)
  (export "main"   (func $main))
  (export "_start" (func $main))
)
\end{minted}

\subsection{Tipos e Valores}

O subconjunto de WebAssembly utilizado neste compilador usa apenas o tipo
$\mathtt{i32}$ (inteiro de 32 bits com e sem sinal, dependendo da instrução).
Todos os valores da IRT, incluindo inteiros, ponteiros de memória e endereços de
registradores temporários, são representados como $\mathtt{i32}$.

\subsection{Gramática do Subconjunto WAT}\label{subsec:wat-grammar}

A figura~\ref{fig:wat-grammar} apresenta a gramática do subconjunto de WAT
gerado pelo compilador.

\begin{figure}[htbp]
\[
\begin{array}{lcll}
  \mathit{mod}  &::=& \mathtt{(module}\ \vec{\mathit{import}}\ \mathit{mem}?\ \vec{\mathit{global}}\ \vec{\mathit{func}}\ \vec{\mathit{export}}\mathtt{)} \\[4pt]
  \mathit{import} &::=& \mathtt{(import}\ s_1\ s_2\ \mathtt{(func}\ \$f\ \mathit{ftype}\mathtt{))} \\
  \mathit{mem}   &::=& \mathtt{(memory}\ n\mathtt{)} \\
  \mathit{global} &::=& \mathtt{(global}\ \$g\ \mathtt{(mut\ i32)}\ \mathtt{(i32.const}\ n\mathtt{))} \\
  \mathit{func}  &::=& \mathtt{(func}\ \$f\ \vec{\mathit{param}}\ \mathit{result}?\ \vec{\mathit{local}}\ \vec{i}\mathtt{)} \\
  \mathit{param} &::=& \mathtt{(param}\ \$x\ \mathtt{i32)} \\
  \mathit{local} &::=& \mathtt{(local}\ \$x\ \mathtt{i32)} \\
  \mathit{result} &::=& \mathtt{(result\ i32)} \\
  \mathit{export} &::=& \mathtt{(export}\ s\ \mathtt{(func}\ \$f\mathtt{))} \\[4pt]
  i &::=& \mathtt{i32.const}\ n \\
    &\mid& \mathtt{local.get}\ \$x \mid \mathtt{local.set}\ \$x \mid \mathtt{local.tee}\ \$x \\
    &\mid& \mathtt{global.get}\ \$g \mid \mathtt{global.set}\ \$g \\
    &\mid& \mathit{binop} \\
    &\mid& \mathtt{i32.eqz} \mid \mathtt{i32.load} \mid \mathtt{i32.store} \\
    &\mid& \mathtt{drop} \mid \mathtt{return} \mid \mathtt{unreachable} \\
    &\mid& \mathtt{call}\ \$f \\
    &\mid& \mathtt{(block}\ \$\ell\ \vec{i}\mathtt{)} \\
    &\mid& \mathtt{(loop}\ \$\ell\ \vec{i}\mathtt{)} \\
    &\mid& \mathtt{(if}\ \mathtt{(then}\ \vec{i}\mathtt{)}\ (\mathtt{(else}\ \vec{i}\mathtt{)})?\ \mathtt{)} \\
    &\mid& \mathtt{br}\ \$\ell \mid \mathtt{br\_if}\ \$\ell \\[4pt]
  \mathit{binop} &::=& \mathtt{i32.add} \mid \mathtt{i32.sub} \mid \mathtt{i32.mul}
                     \mid \mathtt{i32.div\_s} \mid \mathtt{i32.rem\_s} \\
    &\mid& \mathtt{i32.and} \mid \mathtt{i32.or} \mid \mathtt{i32.xor} \\
    &\mid& \mathtt{i32.shl} \mid \mathtt{i32.shr\_s} \mid \mathtt{i32.shr\_u} \\
    &\mid& \mathtt{i32.eq} \mid \mathtt{i32.ne}
         \mid \mathtt{i32.lt\_s} \mid \mathtt{i32.le\_s}
         \mid \mathtt{i32.gt\_s} \mid \mathtt{i32.ge\_s}
\end{array}
\]
\caption{Gramática do subconjunto WAT gerado pelo compilador.}\label{fig:wat-grammar}
\end{figure}

% ============================================================================
\section{Semântica Operacional}\label{sec:wat-semantics}
% ============================================================================

A semântica de WebAssembly é definida sobre uma \textbf{pilha de operandos}.
Instruções consomem valores do topo da pilha e empilham resultados.  O estado
de execução de uma função é um par $(\sigma, \mu)$ onde:

\begin{itemize}
  \item $\sigma = [v_1, \ldots, v_n]$ é a \textbf{pilha de operandos} (o topo é à direita);
  \item $\mu$ é o \textbf{ambiente de variáveis locais} (incluindo parâmetros),
    mapeando nomes a valores $\mathtt{i32}$.
\end{itemize}

A memória linear $\mathbf{M}$ e os globais $\mathbf{G}$ são compartilhados por
todas as funções em execução e são tratados como estado global mutable.

\subsection{Regras de Avaliação de Instruções}

Usamos a notação $(\sigma, \mu) \xrightarrow{i} (\sigma', \mu')$ para
``a instrução $i$ transforma o estado $(\sigma, \mu)$ em $(\sigma', \mu')$''.

\begin{Definition}[Semântica das Instruções Básicas]\label{def:wat-instrs}
\[
\dfrac{}
      {(\sigma, \mu) \xrightarrow{\mathtt{i32.const}\ n} (\sigma \cdot n,\ \mu)}
\quad
\dfrac{\mu(\$x) = v}
      {(\sigma, \mu) \xrightarrow{\mathtt{local.get}\ \$x} (\sigma \cdot v,\ \mu)}
\]

\[
\dfrac{\sigma = \sigma' \cdot v}
      {(\sigma, \mu) \xrightarrow{\mathtt{local.set}\ \$x} (\sigma',\ \mu[\$x \mapsto v])}
\quad
\dfrac{\sigma = \sigma' \cdot v \cdot v}
      {(\sigma, \mu) \xrightarrow{\mathtt{local.tee}\ \$x} (\sigma' \cdot v,\ \mu[\$x \mapsto v])}
\]

\[
\dfrac{\sigma = \sigma' \cdot v_2 \cdot v_1 \quad v = \mathit{eval}(\mathit{op}, v_1, v_2)}
      {(\sigma, \mu) \xrightarrow{\mathit{binop}} (\sigma' \cdot v,\ \mu)}
\]

\[
\dfrac{\sigma = \sigma' \cdot a \quad \mathbf{M}[a] = v}
      {(\sigma, \mu) \xrightarrow{\mathtt{i32.load}} (\sigma' \cdot v,\ \mu)}
\quad
\dfrac{\sigma = \sigma'' \cdot a \cdot v}
      {(\sigma, \mu) \xrightarrow{\mathtt{i32.store}} (\sigma'',\ \mu) \quad \mathbf{M}[a] \leftarrow v}
\]

\[
\dfrac{\sigma = \sigma' \cdot v}
      {(\sigma, \mu) \xrightarrow{\mathtt{drop}} (\sigma',\ \mu)}
\]
\end{Definition}

\subsection{Controle de Fluxo Estruturado}

WebAssembly usa controle de fluxo \textbf{estruturado}: não há \texttt{goto}
incondicional.  Os construtores são:

\begin{itemize}
  \item \texttt{block $\ell$ \{ $\vec{i}$ \}}: executa $\vec{i}$; um
    \texttt{br $\ell$} dentro salta para o fim do bloco (\emph{break}).
  \item \texttt{loop $\ell$ \{ $\vec{i}$ \}}: executa $\vec{i}$; um
    \texttt{br $\ell$} dentro salta para o \emph{início} do laço (\emph{continue}).
  \item \texttt{if \{ $\vec{i}_t$ \} else \{ $\vec{i}_f$ \}}: consome o topo
    da pilha; executa o ramo verdadeiro se $\neq 0$, o falso caso contrário.
  \item \texttt{br\_if $\ell$}: consome o topo da pilha; ramifica se $\neq 0$.
\end{itemize}

O padrão canônico de laço \texttt{while} em WAT combina \texttt{block} e
\texttt{loop}:

\begin{minted}{text}
(block $exit
  (loop $head
    ;; condição de saída: se deve sair, br_if $exit
    <condição>
    br_if $exit
    ;; corpo do laço
    <corpo>
    br $head   ;; volta ao início
  )
)
\end{minted}

% ============================================================================
\section{Tradução da IRT para WAT}\label{sec:irt-to-wat}
% ============================================================================

\subsection{Estrutura Geral do Módulo}

Dado um programa IRT $P = [f_1, \ldots, f_n]$, o módulo WAT gerado é
\emph{auto-contido}: não depende de nenhum módulo externo em tempo de execução.

\[
\mathtt{(module}\ \mathit{WasiImports}\ \mathtt{(memory\ 4)}\ \mathtt{(global\ \$heap\_ptr\ \ldots)}\ \mathit{alloc}\ \llbracket f_1 \rrbracket \cdots \llbracket f_n \rrbracket\ \mathit{print}\ \mathit{read\_int}\ \mathit{Exports}\mathtt{)}
\]

onde:
\begin{itemize}
  \item $\mathit{WasiImports}$ importa diretamente as chamadas de sistema WASI
    \texttt{fd\_write} e \texttt{fd\_read} de \texttt{"wasi\_snapshot\_preview1"};
  \item \texttt{(memory (export "memory") 4)} reserva 4 páginas (256\,KB) de
    memória linear e a exporta (exigido pela interface WASI);
  \item \texttt{\$heap\_ptr} é um global mutável inicializado em \textbf{128};
    os endereços 0–127 são reservados para os buffers de I/O do runtime;
  \item $\mathit{alloc}$ é a função auxiliar de alocação de registros;
  \item $\llbracket f \rrbracket$ é a função compilada de cada definição IRT;
  \item $\mathit{print}$ e $\mathit{read\_int}$ são as funções de I/O do
    runtime, embutidas diretamente no módulo gerado;
  \item $\mathit{Exports}$ inclui \texttt{"main"} e \texttt{"\_start"} (ponto
    de entrada WASI), ambos apontando para \texttt{\$main}.
\end{itemize}

\paragraph{Layout de memória.} O módulo compartilha uma única memória linear
entre o runtime e o programa do usuário:

\begin{center}
\begin{tabular}{cl}
\hline
\textbf{Endereços} & \textbf{Uso} \\
\hline
0–15   & estrutura iovec usada por \texttt{fd\_write}/\texttt{fd\_read} \\
16–31  & slot de retorno \texttt{nwritten}/\texttt{nread} \\
32–63  & buffer de entrada (stdin, 31 bytes úteis + terminador) \\
64–95  & buffer de saída (stdout) \\
96–127 & reservado \\
128+   & heap do programa (\texttt{\$heap\_ptr} inicia aqui) \\
\hline
\end{tabular}
\end{center}

\subsection{Alocador de Memória}\label{subsec:alloc}

O alocador usa a estratégia \emph{bump-pointer}: um ponteiro global
\texttt{\$heap\_ptr} avança a cada alocação.  A função \texttt{alloc(n)}
aloca $n$ palavras de 8 bytes cada e retorna o endereço base:

\[
  \mathit{alloc}(n) \equiv
  \begin{array}[t]{l}
    \mathbf{base} \leftarrow \mathtt{heap\_ptr} \\
    \mathtt{heap\_ptr} \leftarrow \mathtt{heap\_ptr} + n \times 8 \\
    \mathbf{return}\ \mathbf{base}
  \end{array}
\]

Em WAT:

\begin{minted}{text}
(func $alloc (param $n i32) (result i32)
  (local $base i32)
  global.get $heap_ptr
  local.set $base
  global.get $heap_ptr
  local.get $n
  i32.const 8
  i32.mul
  i32.add
  global.set $heap_ptr
  local.get $base
)
\end{minted}

\subsection{Tradução de Definições de Funções}

Seja $f = \mathtt{func}\ \mathit{name}(\vec{p})\ s$ uma definição IRT.  A
função WAT correspondente é:

\[
\llbracket f \rrbracket =
  \mathtt{(func}\ \$\mathit{name}\ \vec{\mathtt{(param}\ \$p_i\ \mathtt{i32)}}\ r\
  \vec{\mathtt{(local}\ \$t_j\ \mathtt{i32)}}\ \llbracket s \rrbracket_{\text{seq}}\
  [\ \mathtt{unreachable}\ ]^{r \neq \varepsilon}\ \mathtt{)}
\]

onde $r = \mathtt{(result\ i32)}$ se o corpo contém algum \texttt{RETURN} com
valor, $r = \varepsilon$ caso contrário, e os $t_j$ são os temporários do corpo
que não são parâmetros.  A instrução \texttt{unreachable} final é necessária
para satisfazer o verificador de tipos WAT quando a função retorna via
\texttt{return} explícito em todos os caminhos de controle.

\subsection{Linearização e Reconstrução de Fluxo de Controle}

A IRT representa sequências com o construtor binário $\mathbf{SEQ}$.  O
primeiro passo da compilação é \textbf{linearizar} o corpo em uma lista plana
de instruções $[s_0, s_1, \ldots, s_{n-1}]$ e construir um \textbf{mapa de
rótulos} $\mathcal{L} : \mathit{Label} \to \mathbb{N}$ mapeando cada rótulo ao
índice em que ele aparece.

\paragraph{Reconhecimento de while.}
Um laço \texttt{while} aparece na sequência linearizada na forma:

\[
  \underbrace{\mathbf{LABEL}\ \ell_h}_{i},\quad
  \underbrace{\mathbf{CJUMP}(e, \ell_a, \ell_b)}_{i+1},\quad
  \underbrace{\text{corpo}}_{\ldots},\quad
  \underbrace{\mathbf{JUMP}(\mathbf{NAME}\ \ell_h)}_{\text{antes de exit}},\quad
  \underbrace{\mathbf{LABEL}\ \ell_{\mathit{exit}}}_{\ldots}
\]

Há duas convenções, dependendo de como o \emph{front-end} gera o \textbf{CJUMP}:

\begin{description}
  \item[Convenção A (TWhile/TImp).] $\ell_a = \ell_{\mathit{body}}$ está
    imediatamente após o CJUMP; $\ell_b = \ell_{\mathit{exit}}$.  A condição
    $e$ é verdadeira quando o laço deve \emph{continuar}.  Geração WAT:
    \[
      \llbracket e \rrbracket \cdot [\mathtt{i32.eqz}] \cdot [\mathtt{br\_if}\ \ell_{\mathit{exit}}]
    \]
  \item[Convenção B (estilo factorial).] $\ell_b = \ell_{\mathit{body}}$ está
    imediatamente após o CJUMP; $\ell_a = \ell_{\mathit{exit}}$.  A condição
    $e$ é verdadeira quando o laço deve \emph{parar}.  Geração WAT:
    \[
      \llbracket e \rrbracket \cdot [\mathtt{br\_if}\ \ell_{\mathit{exit}}]
    \]
\end{description}

Em ambos os casos o código WAT gerado tem a forma:

\begin{minted}{text}
(block $exit
  (loop $head
    <check>          ;; br_if $exit (com ou sem i32.eqz)
    <corpo>
    br $head
  )
)
\end{minted}

\paragraph{Reconhecimento de if/if-else.}
Um condicional aparece na forma:

\[
  \underbrace{\mathbf{CJUMP}(e, \ell_t, \ell_f)}_{i},\quad
  \underbrace{\mathbf{LABEL}\ \ell_t}_{i+1},\quad
  \text{ramo then},\quad
  \underbrace{\mathbf{LABEL}\ \ell_f}_{\ldots}
\]

Para o if-else, o ramo then termina com $\mathbf{JUMP}(\mathbf{NAME}\ \ell_e)$
onde $\ell_e$ é o rótulo de saída após o bloco else.  O código gerado é:

\[
  \llbracket e \rrbracket \cdot [\mathtt{if}\ (\mathtt{then}\ \llbracket\text{then}\rrbracket)\ (\mathtt{else}\ \llbracket\text{else}\rrbracket)]
\]

\subsection{Regras de Tradução de Comandos}\label{subsec:trans-stmts}

As regras $\mathcal{S}\llbracket \cdot \rrbracket$ traduzem comandos IRT não
relacionados ao fluxo de controle (esses são tratados pela reconstrução acima).

\begin{Definition}[Tradução de Comandos]\label{def:trans-stmts}
\begin{align*}
\mathcal{S}\llbracket \mathbf{MOVE}(\mathbf{TEMP}\ t, e) \rrbracket
  &= \mathcal{E}\llbracket e \rrbracket \cdot [\mathtt{local.set}\ \$t] \\
\mathcal{S}\llbracket \mathbf{MOVE}(\mathbf{MEM}(a), e) \rrbracket
  &= \mathcal{E}\llbracket a \rrbracket \cdot \mathcal{E}\llbracket e \rrbracket \cdot [\mathtt{i32.store}] \\
\mathcal{S}\llbracket \mathbf{EXP}(\mathbf{CALL}(\mathbf{NAME}\ f, \vec{e})) \rrbracket
  &= \mathcal{E}\llbracket \vec{e} \rrbracket \cdot [\mathtt{call}\ \$f] \cdot d_f \\
\mathcal{S}\llbracket \mathbf{RETURN}([\ ]) \rrbracket
  &= [\mathtt{return}] \\
\mathcal{S}\llbracket \mathbf{RETURN}([e]) \rrbracket
  &= \mathcal{E}\llbracket e \rrbracket \cdot [\mathtt{return}]
\end{align*}
onde $d_f = [\mathtt{drop}]$ se $f$ retorna um valor, $d_f = [\ ]$ caso contrário.
\end{Definition}

\subsection{Regras de Tradução de Expressões}\label{subsec:trans-exprs}

\begin{Definition}[Tradução de Expressões]\label{def:trans-exprs}
\begin{align*}
\mathcal{E}\llbracket \mathbf{CONST}(n) \rrbracket
  &= [\mathtt{i32.const}\ n] \\
\mathcal{E}\llbracket \mathbf{TEMP}(t) \rrbracket
  &= [\mathtt{local.get}\ \$t] \\
\mathcal{E}\llbracket \mathbf{BINOP}(\mathit{op}, e_1, e_2) \rrbracket
  &= \mathcal{E}\llbracket e_1 \rrbracket \cdot \mathcal{E}\llbracket e_2 \rrbracket \cdot [\mathit{op}_{\mathtt{i32}}] \\
\mathcal{E}\llbracket \mathbf{MEM}(e) \rrbracket
  &= \mathcal{E}\llbracket e \rrbracket \cdot [\mathtt{i32.load}] \\
\mathcal{E}\llbracket \mathbf{CALL}(\mathbf{NAME}\ f, \vec{e}) \rrbracket
  &= \mathcal{E}\llbracket \vec{e} \rrbracket \cdot [\mathtt{call}\ \$f] \\
\mathcal{E}\llbracket \mathbf{ESEQ}(s, e) \rrbracket
  &= \mathcal{S}_0\llbracket s \rrbracket \cdot \mathcal{E}\llbracket e \rrbracket
\end{align*}
onde $\mathcal{S}_0\llbracket \cdot \rrbracket$ é a tradução de comandos
dentro de um ESEQ, que produz efeito nulo na pilha, e
$\mathcal{E}\llbracket \vec{e} \rrbracket = \mathcal{E}\llbracket e_1 \rrbracket \cdots \mathcal{E}\llbracket e_n \rrbracket$.
\end{Definition}

A tabela~\ref{tab:binop} mostra a correspondência entre operadores IRT e
instruções WAT.

\begin{table}[htbp]
\centering
\begin{tabular}{ll|ll|ll}
\hline
\textbf{IRT} & \textbf{WAT} & \textbf{IRT} & \textbf{WAT} & \textbf{IRT} & \textbf{WAT} \\
\hline
$+$  & \texttt{i32.add}   & \&  & \texttt{i32.and}   & $=$  & \texttt{i32.eq}   \\
$-$  & \texttt{i32.sub}   & $|$ & \texttt{i32.or}    & $\neq$ & \texttt{i32.ne} \\
$*$  & \texttt{i32.mul}   & $\oplus$ & \texttt{i32.xor} & $<$ & \texttt{i32.lt\_s} \\
$/s$ & \texttt{i32.div\_s}& $\ll$ & \texttt{i32.shl} & $\leq$ & \texttt{i32.le\_s} \\
$\%s$& \texttt{i32.rem\_s}& $\gg_u$ & \texttt{i32.shr\_u} & $>$ & \texttt{i32.gt\_s} \\
     &                    & $\gg_s$ & \texttt{i32.shr\_s} & $\geq$ & \texttt{i32.ge\_s} \\
\hline
\end{tabular}
\caption{Mapeamento de operadores IRT para instruções WAT.}\label{tab:binop}
\end{table}

% ============================================================================
\section{Implementação em Haskell}\label{sec:wat-haskell}
% ============================================================================

A implementação está organizada em três módulos Haskell dentro do pacote
\texttt{bcc328}:

\begin{center}
\begin{tabular}{ll}
\hline
\textbf{Módulo} & \textbf{Responsabilidade} \\
\hline
\texttt{IR.Backend.WAT.WATSyntax} & Tipos de dados para AST WAT \\
\texttt{IR.Backend.WAT.WATPretty} & Impressão de módulos WAT como texto \\
\texttt{IR.Backend.WAT.IRToWAT}   & Tradução IRT $\to$ WAT \\
\hline
\end{tabular}
\end{center}

\subsection{WATSyntax: Tipos de Dados}

O módulo \texttt{WATSyntax} define a representação interna de um módulo WAT.
Os tipos principais são mostrados a seguir:

\begin{minted}{haskell}
type WIdent = String
data WValType = WI32 deriving (Eq, Show)

data WInstr
  = WI32Const  Int           -- i32.const n
  | WLocalGet  WIdent        -- local.get $x
  | WLocalSet  WIdent        -- local.set $x
  | WI32BinOp  WBinOp        -- i32.<op>
  | WI32Eqz                  -- i32.eqz
  | WI32Load                 -- i32.load
  | WI32Store                -- i32.store
  | WDrop | WReturn | WUnreachable
  | WBlock (Maybe WIdent) [WInstr]
  | WLoop  (Maybe WIdent) [WInstr]
  | WIf    [WInstr] (Maybe [WInstr])
  | WBr WIdent | WBrIf WIdent
  | WCall WIdent
  | WGlobalGet WIdent | WGlobalSet WIdent
  deriving (Eq, Show)

data WFunc = WFunc
  { wfName   :: WIdent
  , wfParams :: [WParam]
  , wfResult :: Maybe WValType
  , wfLocals :: [WLocal]
  , wfBody   :: [WInstr]
  } deriving (Eq, Show)

data WModule = WModule
  { wmImports      :: [WImport]
  , wmMemory       :: Maybe Int
  , wmMemoryExport :: Maybe String  -- nome de export inline da memória
  , wmGlobals      :: [WGlobal]
  , wmFuncs        :: [WFunc]
  , wmExports      :: [WExport]
  , wmRawDecls     :: [String]      -- WAT verbatim (funções de runtime)
  } deriving (Eq, Show)
\end{minted}

\subsection{WATPretty: Impressão}

O módulo \texttt{WATPretty} define uma instância \texttt{Pretty WModule}
usando a biblioteca \texttt{HughesPJ}.  A função central é \texttt{ppInstr}
que converte cada construtor \texttt{WInstr} para o texto WAT correspondente:

\begin{minted}{haskell}
ppInstr :: WInstr -> Doc
ppInstr (WI32Const n)        = text "i32.const" <+> int n
ppInstr (WLocalGet  x)       = text "local.get"  <+> dollar x
ppInstr (WLocalSet  x)       = text "local.set"  <+> dollar x
ppInstr (WI32BinOp op)       = ppBinOp op
ppInstr (WBlock mlbl body)   =
  let hdr = text "block" <> maybe empty (\l -> space <> dollar l) mlbl
  in parens (hdr $$ nest 2 (vcat (map ppInstr body)))
ppInstr (WLoop mlbl body)    =
  let hdr = text "loop" <> maybe empty (\l -> space <> dollar l) mlbl
  in parens (hdr $$ nest 2 (vcat (map ppInstr body)))
ppInstr (WIf thenB mElse)    =
  let thenDoc = parens (text "then" $$ nest 2 (vcat (map ppInstr thenB)))
      elseDoc = maybe empty
                  (\eb -> parens (text "else" $$ nest 2 (vcat (map ppInstr eb))))
                  mElse
  in parens (text "if" $$ nest 2 (thenDoc $$ elseDoc))
\end{minted}

\subsection{IRToWAT: O Gerador de Código}

O módulo \texttt{IRToWAT} é o coração do gerador.  Sua função pública é:

\begin{minted}{haskell}
irToWAT :: Program -> Either String WModule
\end{minted}

O módulo é construído com as seguintes configurações fixas:

\begin{minted}{haskell}
irToWAT prog = do
  let retMap = buildRetMap prog
  funcDefs <- mapM (compileFunc retMap) prog
  return WModule
    { wmImports      = wasiImports          -- fd_write e fd_read
    , wmMemory       = Just 4               -- 4 páginas = 256 KB
    , wmMemoryExport = Just "memory"        -- exigido pela WASI
    , wmGlobals      = [heapPtrGlobal]      -- heap_ptr = 128
    , wmFuncs        = allocFunc : funcDefs
    , wmExports      = [ WExport "main"   "main"
                       , WExport "_start" "main" ]
    , wmRawDecls     = [runtimePrint, runtimeReadInt]
    }
\end{minted}

\noindent onde \texttt{heapPtrGlobal = WGlobal "heap\_ptr" WI32 True 128} e
\texttt{wmRawDecls} contém o WAT verbatim das funções \texttt{\$print} e
\texttt{\$read\_int}, embutidas no módulo sem passar pelo AST \texttt{WInstr}.

A compilação de cada função opera em três fases:

\paragraph{Fase 1: Linearização.}
O corpo (um \texttt{Stmt} com \texttt{SEQ}) é achatado numa lista plana
$[s_0, \ldots, s_{n-1}]$ e convertido num array para acesso $O(1)$ por índice.
Simultaneamente constrói-se o mapa de rótulos $\mathcal{L}$.

\begin{minted}{haskell}
linearize :: Stmt -> [Stmt]
linearize (SEQ s1 s2) = linearize s1 ++ linearize s2
linearize s           = [s]

buildLabelMap :: [Stmt] -> LabelMap
buildLabelMap stmts =
  Map.fromList [(l, i) | (i, LABEL l) <- zip [0..] stmts]
\end{minted}

\paragraph{Fase 2: Coleta de temporários locais.}
Para declarar todos os \texttt{local} necessários, percorre-se o corpo
coletando todos os \texttt{TEMP} e subtraindo os parâmetros:

\begin{minted}{haskell}
collectTempsStmt :: Stmt -> Set Temp
collectTempsExpr :: Expr -> Set Temp
-- ... (percorre toda a AST acumulando temporários)
\end{minted}

\paragraph{Fase 3: Compilação por faixa de índice.}
A função \texttt{compileBlock retMap sv lm from to} compila os comandos no
intervalo semi-aberto $[\mathit{from}, \mathit{to})$.  A recursão examina o
comando no índice \texttt{from} e decide:

\begin{itemize}
  \item Se \texttt{sv[from] = LABEL $\ell_h$} e \texttt{sv[from+1] = CJUMP
    ...}, tenta \texttt{tryWhile} para reconhecer o padrão de laço.  Se
    bem-sucedido, gera \texttt{WBlock}/\texttt{WLoop} e compila o resto
    recursivamente.
  \item Se \texttt{sv[from] = CJUMP} com rótulos válidos, reconhece um
    condicional e gera \texttt{WIf}.
  \item \texttt{JUMP} é silenciosamente ignorado (já consumido nos padrões
    acima ou é um desvio redundante).
  \item Qualquer outro comando vai para \texttt{compileSingleStmt}.
\end{itemize}

\begin{minted}{haskell}
tryWhile :: LabelMap -> SArr -> Int -> Label -> Label -> Label
         -> Maybe (Int, Int, Label, Bool)
tryWhile lm sv fallThrough la lb lh =
  tryConvA <|> tryConvB
  where
    -- Convenção A: la=corpo (logo após CJUMP), lb=saída
    tryConvA = do
      posLa <- Map.lookup la lm
      guard (posLa == fallThrough)
      posLb <- Map.lookup lb lm
      guard (posLb > fallThrough)
      JUMP (NAME lh') <- Just (sv ! (posLb - 1))
      guard (lh' == lh)
      return (posLa, posLb, lb, False)  -- exitCond=False: usar i32.eqz

    -- Convenção B: lb=corpo (logo após CJUMP), la=saída
    tryConvB = do
      posLb <- Map.lookup lb lm
      guard (posLb == fallThrough)
      posLa <- Map.lookup la lm
      guard (posLa > fallThrough)
      JUMP (NAME lh') <- Just (sv ! (posLa - 1))
      guard (lh' == lh)
      return (posLb, posLa, la, True)   -- exitCond=True: br_if direto
\end{minted}

% ============================================================================
\section{O Runtime de Entrada e Saída}\label{sec:wat-runtime}
% ============================================================================

As funções \texttt{\$print} e \texttt{\$read\_int} são embutidas diretamente no
módulo WAT gerado, usando as chamadas de sistema da interface WASI
(\emph{WebAssembly System Interface}).  Não há módulo de runtime separado — o
código gerado pelo compilador já é um módulo executável completo.

\subsection{Funções WASI utilizadas}

O compilador importa duas primitivas WASI:

\begin{center}
\begin{tabular}{lll}
\hline
\textbf{Função WASI} & \textbf{Assinatura} & \textbf{Uso} \\
\hline
\texttt{fd\_write} & \texttt{(i32 i32 i32 i32) → i32} & escrever em stdout \\
\texttt{fd\_read}  & \texttt{(i32 i32 i32 i32) → i32} & ler de stdin \\
\hline
\end{tabular}
\end{center}

Ambas recebem: \texttt{(fd, iovs\_ptr, iovs\_len, result\_ptr)}, onde
\texttt{iovs\_ptr} aponta para uma estrutura iovec armazenada na memória linear.

\subsection{Função \texttt{\$print}}

Converte um \texttt{i32} para decimal, escreve os dígitos no buffer de saída
(endereços 64–95) e chama \texttt{fd\_write}:

\begin{minted}{text}
(func $print (param $value i32)
  (local $temp i32) (local $digit i32)
  (local $len  i32) (local $neg   i32) (local $pos  i32)
  ;; abs($value) → $temp; flag de sinal → $neg
  ;; constrói string decimal de trás para frente a partir do endereço 95
  ;; prepende '-' se negativo; acrescenta '\n'
  ;; configura iovec[0] = {ptr, len} e chama fd_write(1, 0, 1, 16)
)
\end{minted}

\subsection{Função \texttt{\$read\_int}}

Configura a iovec para ler até 31 bytes de stdin no buffer de entrada (endereços
32–62), depois percorre o buffer convertendo dígitos ASCII em \texttt{i32}:

\begin{minted}{text}
(func $read_int (result i32)
  (local $result i32) (local $pos i32)
  (local $char   i32) (local $neg i32)
  ;; configura iovec[0] = {32, 31} e chama fd_read(0, 0, 1, 16)
  ;; detecta sinal '-'; percorre dígitos com i32.load8_u
  ;; retorna $result (ou -$result se negativo)
)
\end{minted}

O código WAT completo de ambas as funções está no campo \texttt{wmRawDecls} do
módulo gerado e pode ser inspecionado em
\texttt{workspace/src/IR/Backend/WAT/IRToWAT.hs}.

% ============================================================================
\section{Execução de Programas WebAssembly}\label{sec:wat-exec}
% ============================================================================

O flake deste projeto disponibiliza as ferramentas \texttt{wabt} (que inclui
\texttt{wat2wasm} e \texttt{wasm-validate}) e \texttt{wasmtime}.

\subsection{Pipeline Completo}

Como o módulo WAT gerado é auto-contido (inclui as funções de I/O e importa
diretamente a interface WASI), o pipeline de compilação tem apenas três passos:

\begin{enumerate}
  \item \textbf{Compilar TImp para IRT:}
\begin{minted}{bash}
$ cabal run timp -- --ir -f prog.timp
\end{minted}
    Gera \texttt{prog.ir}.

  \item \textbf{Compilar IRT para WAT:}
\begin{minted}{bash}
$ cabal run ir -- --wat prog.ir
\end{minted}
    Gera \texttt{prog.wat}, que já contém as funções \texttt{\$print} e
    \texttt{\$read\_int} embutidas e importa \texttt{fd\_write}/\texttt{fd\_read}
    diretamente de \texttt{wasi\_snapshot\_preview1}.

  \item \textbf{Montar e executar:}
\begin{minted}{bash}
$ wat2wasm prog.wat -o prog.wasm
$ echo "5" | wasmtime prog.wasm
\end{minted}
\end{enumerate}

Não é necessário compilar nem mesclar nenhum módulo de runtime separado.

\subsection{Exemplo: Fatorial}\label{subsec:wat-factorial}

Considere o programa IRT \texttt{factorial.ir} (gerado a partir de TImp):

\begin{minted}{text}
func main() {
  t0 := call read_int();
  t1 := call fact(t0);
  call print(t1);
  return;
}

func fact(n) {
  t0 := 1;
L_head:
  cjump n <= 0 L_done L_body;
L_body:
  t0 := t0 * n;
  n  := n - 1;
  jump L_head;
L_done:
  return t0;
}
\end{minted}

O gerador produz (com nomes simplificados para legibilidade):

\begin{minted}{text}
(func $fact (param $n i32) (result i32)
  (local $t0 i32)
  i32.const 1
  local.set $t0
  (block $L_done
    (loop $L_head
      local.get $n
      i32.const 0
      i32.le_s
      br_if $L_done         ;; Convenção B: condição de saída direta
      local.get $t0
      local.get $n
      i32.mul
      local.set $t0
      local.get $n
      i32.const 1
      i32.sub
      local.set $n
      br $L_head
    )
  )
  local.get $t0
  return
  unreachable
)
\end{minted}

\subsection{Verificação com wasm-validate}

Antes de executar, é útil validar o módulo gerado:

\begin{minted}{bash}
$ wasm-validate prog.wasm
\end{minted}

Se não houver saída, o módulo é válido.  Erros de tipo, como o caso de uma
função com \texttt{(result i32)} que não deixa valor na pilha, são detectados aqui.

% ============================================================================
% ============================================================================
\section{Conclusão}\label{sec:wat-conclusao}
% ============================================================================

Este capítulo fechou o \emph{pipeline} de compilação das linguagens TWhile e
TImp, traduzindo a representação intermediária em árvore (IRT) para módulos
WebAssembly Text (WAT) executáveis.

Para cada etapa foram apresentados:
\begin{itemize}
  \item a definição formal do subconjunto WAT gerado (gramática e semântica
    operacional da máquina de pilha);
  \item as regras de tradução dirigida por sintaxe
    $\mathcal{E}\llbracket \cdot \rrbracket$ e
    $\mathcal{S}\llbracket \cdot \rrbracket$, que cobrem expressões,
    comandos, chamadas e alocação de memória;
  \item o algoritmo de reconstrução de fluxo de controle estruturado a partir
    da IRT linearizada, com suporte às duas convenções de geração de CJUMP;
  \item a implementação em Haskell nos módulos \texttt{WATSyntax},
    \texttt{WATPretty} e \texttt{IRToWAT}, evidenciando a correspondência
    direta entre teoria e código;
  \item o runtime WASI embutido no módulo gerado (funções \texttt{\$print} e
    \texttt{\$read\_int} via \texttt{wmRawDecls}) e o pipeline de execução
    completo em três passos: \texttt{timp --ir}, \texttt{ir --wat},
    \texttt{wat2wasm} + \texttt{wasmtime}.
\end{itemize}

A escolha de WebAssembly como alvo final permite executar programas compilados
em qualquer ambiente que suporte WASI, como navegadores, servidores ou contêineres,
sem dependência de uma arquitetura de hardware específica.  Isso encerra, com
uma implementação concreta e executável, o ciclo completo do compilador
apresentado na disciplina: da análise léxica e sintática, passando pela
verificação de tipos, geração e otimização da IRT, até a produção de código
executável.
